\subsection{机制：不对称的假翻转抑制}\label{sec:mechanism}

第~\ref{sec:formal} 与第~\ref{sec:temporal} 小节把窗口定位在时间上；这里我们考察偏移对它触及的那一轮做了什么。一个配对探针在 RAND 下译码 20 把密钥、每把 1{,}000 个错误，得到 20{,}000 条首轮记录。平均首轮总翻转从 C0 的 29.00 降到 FRO-2 的 8.74。错误翻转下降 $9.5\times$（约为原来的十分之一），正确翻转只下降 $2.7\times$（约为原来的 38\%）。首轮后的残余错误重量从 79.02 升到 84.94。因此偏移用首轮未做完的工作，换取对假翻转的选择性。

失败走一般的译码路径。停滞占 RAND 下 C0 失败的 198/282，占 TYPE-I 的 257/372，占 NCW-shallow 的 101/141。归档分类器与预指定的四分之一预算窗口合并后没有出入，因为每次失败都用尽迭代预算。这些计数使用真实错误支持集，因此它们诊断机制，而不是给译码器一个它能自己算出的统计量。抬高阈值削掉的错误翻转远多于正确翻转；代价是正确翻转被推迟。

\dblfig{fig3_mechanism.pdf}{首轮机制。(a)~错误翻转，(b)~正确翻转，(c)~首轮后残余错误重量，单位为每次译码的平均位置数。柱为 20 把密钥、每臂 1{,}000 次译码的归档均值；每对柱上方的标签给出 FRO-2/C0。三个量使用各自独立的纵轴。}{fig:mechanism}

\subsection{迁移与稳健性}\label{sec:transfer}

把一个窗口固定在一个工作点之后，我们检验它能否经受工作点、重量、预算与语义的改变。在 $P_2=(3001,37,66)$，调度无需重新调参即完成迁移，覆盖 30 把密钥、每把 3{,}000 个错误，每变体 90{,}000 次译码。RR 在 RAND 下为 1.95 [1.58, 2.42]，TYPE-I 下为 3.21 [2.57, 4.14]，NCW-shallow 下为 2.51 [1.85, 3.46]；DFR 分别从 $2.278\times10^{-3}$ 降到 $1.167\times10^{-3}$、从 $3.000\times10^{-3}$ 降到 $9.333\times10^{-4}$、从 $1.200\times10^{-3}$ 降到 $4.778\times10^{-4}$。最差组改善为 0.384。这次运行与扩大确认是同一名义点上的两场独立战役（30 把密钥各 3{,}000 个错误，对 50 把密钥各 1{,}000 个），所以两者的 RR 不必重合。图~\ref{fig:scope}(a) 展示这三行。

次级检查保持同号。$t=94$ 与 $t=96$ 处 RR 为 2.82 [2.28, 3.60] 与 1.87 [1.72, 2.04]；最大块 $(6007,56,100)$ 给出 2.41 [2.09, 2.82]，从 $8.125\times10^{-3}$ 降到 $3.375\times10^{-3}$；预算 7、10、20 给出 1.72--3.17。12 个设置中 6 个低于可测带、6 个高于可测带。低于可测带的 5 个双方均无失败，第 6 个为 FRO-2 1/30{,}000、C0 0/30{,}000，因此低率侧无法确立方向。高于可测带的 4 个双方均有失败且 C0 失败更多（RR 1.77、1.08、1.012 与 1.004），另外 2 个在 30{,}000/30{,}000 处饱和，没有区分力。扫描在能分辨处与该方向一致，但并不能在每一个设置上确立该方向。动态距离实现在 RAND 与 TYPE-I 下复现方向（1.59 与 2.01；$z=7.94$ 与 $11.98$），NCW 未分辨（$z=1.36$）。

最低端点把范围向下扩展：C0 在 1{,}840{,}000 次译码中失败 130 次（$7.07\times10^{-5}$），FRO-2 失败 35 次，RR 3.714 [2.69, 5.57]。在留出点 $P_4=(3607,39,76)$，一项预指定规则选中 $\delta=1$；验证记录 C0 失败 40 次，而每个一次性偏移各失败 8 次（RR 5.00）。可测设置保持方向；幅度变化，低失败率设置无法分辨。

\dblfig{fig5_scope_and_boundary.pdf}{作用范围与边界。(a)~独立 P2 迁移：C0/FRO-2 风险比与 95\% 归档区间。(b)~构造的近码字交叠 $\kappa=0,12,46,60$ 处的绝对 DFR；标记的水平偏移只为避免遮挡。(c)~同一组交叠上的 C0/FRO-2 风险比；$\kappa=60$ 时两种方法均饱和，因此不画比值。}{fig:scope}
